\documentclass[11pt,a4paper,reqno]{amsart}

\usepackage[utf8]{inputenc}
\usepackage[T1]{fontenc}
\usepackage{amsmath,amssymb,amsfonts,amsthm,mathtools}
\usepackage{geometry}
\geometry{margin=2.5cm}
\usepackage[dvipsnames]{xcolor}
\usepackage{hyperref}
\hypersetup{
    colorlinks=true,
    linkcolor=blue!70!black,
    citecolor=green!50!black,
    urlcolor=blue!70!black
}
\usepackage{microtype}
\usepackage{booktabs}
\usepackage{array}
\usepackage{physics}

% Theorem Environments
\newtheorem{theorem}{Theorem}[section]
\newtheorem{lemma}[theorem]{Lemma}
\newtheorem{proposition}[theorem]{Proposition}
\newtheorem{corollary}[theorem]{Corollary}

\theoremstyle{definition}
\newtheorem{definition}[theorem]{Definition}
\newtheorem{example}[theorem]{Example}

\theoremstyle{remark}
\newtheorem{remark}[theorem]{Remark}

\numberwithin{equation}{section}
\raggedbottom

% Custom Operators
\newcommand{\dif}{\mathrm{d}}
\DeclareMathOperator{\Sym}{Sym}
\DeclareMathOperator{\End}{End}
\DeclareMathOperator{\supp}{supp}
\DeclareMathOperator{\im}{im}
\DeclareMathOperator{\Cov}{Cov}
\DeclareMathOperator{\Var}{Var}
\DeclareMathOperator{\Unp}{Unp}
\DeclareMathOperator{\Proj}{Proj}
\DeclareMathOperator{\reach}{reach}

\title[Beyond the Spectrum III]{Beyond the Spectrum III: Higher Invariants, Bipartite Kernels, and Non-Equilibrium Geometry of Functional Realizations}

\author{Reinaldo M. Silva-Filho}
\address{Programa de P\'os-Gradua\c{c}\~ao em Engenharia de Sistemas e Automa\c{c}\~ao (PPGEEAA), Departamento de Engenharia de Sistemas (DES), Universidade Federal de Lavras (UFLA), Lavras, MG, Brazil}
\email{reinaldo.silva@ufla.br}

\date{\today}

\begin{document}

\begin{abstract}
We establish the third installment of the functional realization framework for matrices and higher-order tensors, bridging continuous linear algebra with symplectic topology, microlocal sheaf theory, nonlinear spectral geometry, non-equilibrium stochastic thermodynamics, holographic quantum information, and geometric measure theory. 
First, we construct the generalized Ekeland--Hofer capacities $c_k^{\mathrm{EH}}(A, t)$ and Hofer--Zehnder invariants $c_{\mathrm{HZ}}(A, t)$ of functional sublevel sets $K_t(A) \subset \mathbb{R}^{2n}$, proving sharp $C^0$-Lipschitz stability for Viterbo--Floer spectral invariants under quadratic convex growth. 
Second, we advance the codomain of functional realizations to matrix-valued bundles $\Phi(A) \in C^\infty(\Omega, \mathrm{Sym}_n(\mathbb{R}))$, proving that transversal rank stratifications $\Omega = \bigsqcup \Sigma_r(A)$ satisfy Whitney conditions (A) and (B) via Mather's stratified transversality theorem, while the characteristic cycle $CC(\mathcal{F}_A)$ of constructible vanishing cycle complexes computes global Euler characteristics via Kashiwara's index theorem. 
Third, we analyze the nonlinear Rayleigh--Finsler variational spectrum of the weighted $p$-Laplacian $\Delta_p^A$ on connected domains, rigorously establishing the Kawohl--Fridman Cheeger limit $\lim_{p \to \infty} (\lambda_1^{(p)}(A))^{1/p} = h(A)$ through two-sided variational bounds and bounding the Gromov $\delta$-hyperbolicity of complete conformal tensor metrics. 
Fourth, we formulate overdamped Langevin diffusions under the matrix potential $V_A = -\log \Phi(A)$, establishing an exact Jarzynski free energy identity for unnormalized tensor protocols where $\Delta F(A, B) = -\log(\Tr(B)/\Tr(A))$, and deriving finite-time dissipation bounds where the thermodynamic length $\mathcal{L}(A, B)$ lower-bounds the quadratic Wasserstein distance $\mathcal{W}_2$. 
Fifth, we introduce the bipartite kernel operator $K_A(x, y) \in L^2(\Omega \times \Omega, \mathbb{C})$, constructing continuous modular Hamiltonians on non-degenerate support subspaces $\mathcal{H}_{\mathrm{supp}} = (\ker \rho_{\Omega_1})^\perp$, reflected entropies $S_R(\Omega_1 : \Omega_2)$, and proving the holographic entanglement wedge cross-section inequality $S_R \ge 2 E_W$. 
Finally, we construct simplicial Mellin transforms $\mathcal{Z}_A(s)$ over the continuous Pascal simplex $\Delta_m$, extracting the discrete Barnes--Kigami residue spectrum, proving the exact fractal spectral dimension identity $d_s = \frac{2\log N}{\log(N/\rho)}$ (verified on the Sierpinski gasket benchmark), and proving quantitative lower bounds on the Federer reach $\mathrm{reach}(\Sigma_t)$ incorporating both curvature and non-local separation. 
All 21 theoretical results are formally catalogued in a dual obligation ledger and aligned with machine-checked Lean~4 specifications.
\end{abstract}

\maketitle
\tableofcontents

% =========================================================================
% SECTION 1
% =========================================================================
\section{Introduction and the Higher Categorical Frontier}
\label{sec:intro}

The classical spectral theory of matrices and tensors characterizes an algebraic object $A \in \mathcal{S}_+^n$ through scalar invariants invariant under orthogonal or unitary transformations: eigenvalues $\lambda_j(A)$, traces $\Tr(A^k)$, and determinants $\det(A)$. In Volumes~I and~II of this series, we introduced the paradigm of \emph{functional realizations}, embedding finite- and infinite-dimensional tensors into continuous function spaces:
\begin{equation}
\Phi \colon \mathcal{T}^k(V) \longrightarrow \mathcal{F}(\Omega, \mathbb{R}), \qquad A \longmapsto \Phi(A)(x),
\end{equation}
where $\Omega \subset \mathbb{R}^d$ is a smooth compact domain or Riemannian manifold. This continuous lifting unlocked the rich geometric measures of Sobolev spaces $W^{1,p}$, bounded variation $\mathrm{BV}(\Omega)$, continuous optimal transport $\mathcal{W}_2(\mu_A, \mu_B)$, Bakry--Émery Ricci curvature lower bounds $\kappa_{\mathrm{BE}}(A)$, persistent homology diagrams $\mathrm{Dgm}_k(A)$, integrated Willmore bending energies $\mathcal{W}(\Phi(A))$, and non-commutative Dixmier traces $\Tr_\omega(\Phi(A) |\mathcal{D}|^{-d})$.

Despite the analytical breadth of these invariants, Volumes~I and~II operated under two fundamental structural constraints:
\begin{enumerate}
    \item \textbf{Scalar Codomain Restriction:} The realization $\Phi(A)$ was strictly a scalar field $\Omega \to \mathbb{R}$ or a Radon probability measure $\dif\mu_A = \Phi(A) \dif x$. Consequently, topological phase-space rigidity (symplectic capacity), gauge-field parallel transport, and internal quantum entanglement remained inaccessible.
    \item \textbf{Static Equilibrium Assumption:} Invariants were computed on static tensors or quasistatic limits, blind to finite-time non-equilibrium dissipation, dynamical fluctuations, and irreversible entropy production.
\end{enumerate}

The present monograph (\emph{Beyond the Spectrum III}) breaks through these boundaries along three methodological fronts:
\begin{itemize}
    \item \textbf{Symplectic and Contact Invariance:} In Section~\ref{sec:symplectic}, we treat $\Omega = \mathbb{R}^{2n}$ as a symplectic phase space $(\Omega, \omega_0)$. Rather than volume or $L^p$ norms, we compute the Ekeland--Hofer capacities $c_k^{\mathrm{EH}}(A, t)$ and Hofer--Zehnder invariants $c_{\mathrm{HZ}}(A, t)$ of the sublevel sets $K_t(A)$, capturing the non-squeezing geometry of Hamilton's equations.
    \item \textbf{Codomain Promotion to Matrix Bundles and Bipartite Kernels:} In Section~\ref{sec:microlocal}, $\Phi(A)(x)$ is promoted to a smooth matrix-valued field in $\mathrm{Sym}_n(\mathbb{R})$, unlocking Whitney rank stratifications and constructible characteristic cycles. In Section~\ref{sec:bipartite}, we perform the fundamental codomain leap to bipartite integral kernels $K_A(x, y) \in L^2(\Omega \times \Omega, \mathbb{C})$, constructing continuous modular Hamiltonians on their support subspaces, canonical GNS purifications, reflected entropies, and proving holographic entanglement wedge cross-section inequalities.
    \item \textbf{Nonlinear Analysis, Stochastic Thermodynamics, and Fractals:} In Section~\ref{sec:nonlinear}, we investigate the continuous $p$-Laplacian family $\{\lambda_1^{(p)}(A)\}_{p \in (1,\infty)}$ on connected domains and Gromov $\delta$-hyperbolicity. In Section~\ref{sec:thermo}, we model Langevin diffusions under the matrix potential $V_A = -\log \Phi(A)$, establishing Jarzynski's identity for unnormalized tensors and finite-time thermodynamic length bounds. Finally, Sections~\ref{sec:simplicial} and~\ref{sec:federer} develop the continuous Pascal simplex Mellin--Barnes residue spectrum and Federer reach geometry.
\end{itemize}

Every definition and theorem in this treatise is assigned an immutable obligation tag ($\mathtt{OBL\text{-}001}$ to $\mathtt{OBL\text{-}021}$), synchronized 1:1 with an accompanying formal specification package written in the Lean~4 interactive theorem prover.

% =========================================================================
% SECTION 2
% =========================================================================
\section{Symplectic and Floer Invariants of Functional Sublevel Sets}
\label{sec:symplectic}

Let $(\mathbb{R}^{2n}, \omega_0)$ be the standard symplectic vector space with canonical coordinates $x = (q_1, \dots, q_n, p_1, \dots, p_n)$ and non-degenerate closed 2-form $\omega_0 = \sum_{j=1}^n \dif p_j \wedge \dif q_j$. 
Let $A \in \mathcal{S}_+^n$ be a positive definite matrix, and let $\Phi(A) \in C^2(\mathbb{R}^{2n}, \mathbb{R})$ be a strictly convex functional realization satisfying $\nabla^2 \Phi(A)(x) \ge \mu I > 0$ for all $x \in \mathbb{R}^{2n}$, with $\lim_{\|x\| \to \infty} \Phi(A)(x) = +\infty$.

For any regular value $t > \min_{x} \Phi(A)(x)$, the sublevel set
\begin{equation}
K_t(A) \coloneqq \left\{ x \in \mathbb{R}^{2n} \mid \Phi(A)(x) \le t \right\}
\end{equation}
is a compact, strictly convex body in $\mathbb{R}^{2n}$ with smooth boundary $\Sigma_t \coloneqq \partial K_t(A)$. On this compact body, we define the sharp local Hessian spectral bounds:
\begin{equation}
\lambda_{\max} \coloneqq \sup_{x \in K_t(A)} \lambda_{\max}\left(\nabla^2 \Phi(A)(x)\right) < \infty, \qquad \lambda_{\min} \coloneqq \inf_{x \in K_t(A)} \lambda_{\min}\left(\nabla^2 \Phi(A)(x)\right) \ge \mu > 0.
\end{equation}

\begin{theorem}[\textbf{OBL-001: Sublevel Convexity and Ekeland--Hofer Capacity}]
\label{thm:eh_capacity_pos}
Let $A \in \mathcal{S}_+^n$ and let $\Phi(A) \in C^2(\mathbb{R}^{2n})$ be strictly convex with $\nabla^2 \Phi(A) \ge \mu I > 0$. For any regular value $t > \min \Phi(A)$, the first Ekeland--Hofer capacity $c_1^{\mathrm{EH}}(K_t(A))$ is strictly positive and satisfies the two-sided variational bound:
\begin{equation}
\label{eq:eh_bound}
0 < \frac{2\pi (t - \min \Phi(A))}{\lambda_{\max}} \le c_1^{\mathrm{EH}}(K_t(A)) \le \frac{2\pi (t - \min \Phi(A))}{\lambda_{\min}} < \infty.
\end{equation}
\end{theorem}

\begin{proof}
Under strict convexity, the Clarke dual action functional on the loop space $H^{1/2}(S^1, \mathbb{R}^{2n})$ is given by
\begin{equation}
\Psi(z) = \frac{1}{2} \int_0^1 \langle -J_0 \dot{z}, z \rangle \dif \tau + \int_0^1 H^*(\dot{z}) \dif \tau,
\end{equation}
where $H^*$ is the Legendre--Fenchel transform of the shifted Hamiltonian $H(x) = (\Phi(A)(x) - \min \Phi(A)) / (t - \min \Phi(A))$. By the Ekeland--Hofer minimax principle \cite{EkelandHofer1989}, $c_1^{\mathrm{EH}}(K_t(A))$ equals the critical value of $\Psi$ corresponding to the first non-trivial Clarke dual action. Because $\lambda_{\min} I \le \nabla^2 \Phi(A)(x) \le \lambda_{\max} I$ for all $x \in K_t(A)$, Taylor expansion centered at the unique minimum $x_0$ yields the geometric inclusion:
\begin{equation}
E_{\mathrm{in}} \subseteq K_t(A) \subseteq E_{\mathrm{out}},
\end{equation}
where $E_{\mathrm{in}}$ is an inscribed ellipsoid of minimum semi-axis $r_{\min} = \sqrt{2(t - \min \Phi(A)) / \lambda_{\max}}$ and $E_{\mathrm{out}}$ is a circumscribed ellipsoid of maximum semi-axis $R_{\max} = \sqrt{2(t - \min \Phi(A)) / \lambda_{\min}}$. By the monotonicity and conformality axioms of symplectic capacities, $c_1^{\mathrm{EH}}(E_{\mathrm{in}}) \le c_1^{\mathrm{EH}}(K_t(A)) \le c_1^{\mathrm{EH}}(E_{\mathrm{out}})$. Since the capacity of an ellipsoid is $\pi r_{\min}^2$, we obtain $\pi r_{\min}^2 = \frac{2\pi(t - \min \Phi(A))}{\lambda_{\max}}$ and $\pi R_{\max}^2 = \frac{2\pi(t - \min \Phi(A))}{\lambda_{\min}}$, establishing \eqref{eq:eh_bound}.
\end{proof}

\begin{proposition}[\textbf{OBL-002: Hofer--Zehnder Periodic Orbit Invariance}]
\label{prop:hz_symp_inv}
Let $\psi \in \mathrm{Symp}(\mathbb{R}^{2n}, \omega_0)$ be any smooth symplectomorphism ($\psi^* \omega_0 = \omega_0$). Then the Hofer--Zehnder capacity $c_{\mathrm{HZ}}$ of the functional domain is invariant:
\begin{equation}
c_{\mathrm{HZ}}(\psi(K_t(A))) = c_{\mathrm{HZ}}(K_t(A)) = \inf_{\gamma \subset \partial K_t(A)} \left| \oint_\gamma p \dif q \right|,
\end{equation}
where the infimum runs over all closed characteristic curves of the Hamiltonian vector field $X_{\Phi(A)}$ on $\partial K_t(A)$.
\end{proposition}

\begin{proof}
For strictly convex bodies with smooth boundary in $\mathbb{R}^{2n}$, the foundational capacity theory of Hofer and Zehnder \cite{HoferZehnder1994} establishes that $c_{\mathrm{HZ}}(K_t(A))$ coincides with the minimal action of a closed characteristic on $\partial K_t(A)$. Under a symplectomorphism $\psi$, the Hamiltonian flow generated by $H \circ \psi^{-1}$ on $\psi(K_t(A))$ is conjugate to the flow of $H$ on $K_t(A)$ via $\psi$. Symplectic action $\oint_\gamma \lambda$ is invariant under symplectomorphisms since $\psi^* \lambda = \lambda + \dif f$. Thus, the action spectrum $\mathcal{S}(\partial K_t(A))$ is preserved, proving invariance.
\end{proof}

\begin{theorem}[\textbf{OBL-003: Floer--Viterbo Spectral Lipschitz Stability}]
\label{thm:viterbo_lipschitz}
Let $\Phi(A), \Phi(B) \in C^2(\mathbb{R}^{2n})$ have uniform quadratic convex growth at infinity: $c_1 \|x\|^2 \le \Phi(A)(x), \Phi(B)(x) \le c_2 \|x\|^2$ for $\|x\| \ge R$. Let $a \in H_*(\mathbb{R}^{2n}; \mathbb{Z}_2) \setminus \{0\}$, and let $K \subset \mathbb{R}^{2n}$ be any compact domain containing all periodic orbits of action below the spectral cutoff. The Viterbo spectral invariant $c(a, \Phi(A))$ satisfies:
\begin{equation}
\label{eq:viterbo_lip}
\left| c(a, \Phi(A)) - c(a, \Phi(B)) \right| \le \|\Phi(A) - \Phi(B)\|_{L^\infty(K)}.
\end{equation}
\end{theorem}

\begin{proof}
By Viterbo's generating function theory \cite{Viterbo1992} and filtered Hamiltonian Floer homology for asymptotically quadratic Hamiltonians, the spectral invariant $c(a, H)$ is constructed via a minimax critical value on the loop space. The filtered Floer continuation morphism $\Phi_{\mathrm{cont}} \colon CF_*^{\lambda}(H_A) \to CF_*^{\lambda + \Delta}(H_B)$ is well-defined, where the action shift satisfies $\Delta \le \int_0^1 \max_{x \in K} (H_B - H_A) \dif t \le \|H_B - H_A\|_{L^\infty(K)}$. Monotonicity of the spectral selector gives $c(a, H_B) \le c(a, H_A) + \|H_B - H_A\|_{L^\infty(K)}$, and swapping indices yields \eqref{eq:viterbo_lip}.
\end{proof}

% =========================================================================
% SECTION 3
% =========================================================================
\section{Microlocal Sheaves and Stratified Rank Characteristic Cycles}
\label{sec:microlocal}

We now perform the structural promotion of the realization map to smooth matrix-valued fields $\Phi(A) \in C^\infty(\Omega, \mathrm{Sym}_n(\mathbb{R}))$.

\begin{definition}[\textbf{OBL-004: Rank Stratification}]
\label{def:rank_strat}
Let $\Phi(A) \in C^\infty(\Omega, \mathrm{Sym}_n(\mathbb{R}))$. For each integer $r \in \{0, 1, \dots, n\}$, the rank-$r$ stratum of $\Omega$ is defined by:
\begin{equation}
\Sigma_r(A) \coloneqq \left\{ x \in \Omega \mid \operatorname{rank}(\Phi(A)(x)) = r \right\}.
\end{equation}
\end{definition}

\begin{proposition}[\textbf{OBL-004: Whitney Regularity of Rank Strata}]
\label{prop:whitney_rank_strat}
Assume that the mapping $\Phi(A) \colon \Omega \to \mathrm{Sym}_n(\mathbb{R})$ is stratified transverse (satisfying Mather's condition $(a_f)$) to the natural rank stratification $\mathcal{R} = \{\mathcal{R}_r\}_{r=0}^n$ of $\mathrm{Sym}_n(\mathbb{R})$. Then $\Omega = \bigsqcup_{r=0}^n \Sigma_r(A)$ is a regular Whitney stratification satisfying conditions~(A) and~(B).
\end{proposition}

\begin{proof}
The space $\mathrm{Sym}_n(\mathbb{R})$ admits a canonical Whitney stratification by the algebraic submanifolds $\mathcal{R}_r$ of symmetric matrices of exact rank $r$, where each stratum has codimension $\binom{n-r+1}{2}$. By Mather's stratified transversality theorem, if a smooth map is transverse to a Whitney stratified space in the sense of Thom--Mather, the preimage stratification $\Sigma_r(A) = \Phi(A)^{-1}(\mathcal{R}_r)$ automatically inherits Whitney's conditions (A) and (B), and local triviality holds along each stratum.
\end{proof}

\begin{lemma}[\textbf{OBL-005: Kashiwara--Schapira Micro-Support Lagrangian Property}]
\label{lem:microsupport_lagrangian}
Let $\mathcal{F}_A \in D_c^b(\Omega)$ be the constructible sheaf complex of vanishing cycles of $\det \Phi(A)$, which is locally constant along each stratum $\Sigma_r(A)$. Then the micro-support $SS(\mathcal{F}_A) \subset T^*\Omega$ is a closed conic Lagrangian subvariety, given by:
\begin{equation}
SS(\mathcal{F}_A) \subseteq \bigcup_{r=0}^n \overline{T^*_{\Sigma_r(A)} \Omega}.
\end{equation}
\end{lemma}

\begin{proof}
Because $\mathcal{F}_A$ is constructible with respect to the Whitney stratification $\Sigma_\bullet(A)$, Kashiwara--Schapira involutivity \cite{KashiwaraSchapira1990} guarantees that $SS(\mathcal{F}_A)$ is isotropic and coisotropic in the canonical symplectic manifold $(T^*\Omega, \dif p \wedge \dif q)$, hence purely Lagrangian. Since $\mathcal{F}_A$ has locally constant cohomology sheaves on each stratum $\Sigma_r(A)$, the micro-support is contained in the union of the closures of the conormal bundles $\overline{T^*_{\Sigma_r(A)} \Omega}$.
\end{proof}

\begin{theorem}[\textbf{OBL-006: Characteristic Cycle Kashiwara Index Identity}]
\label{thm:kashiwara_index_cc}
Let $\Omega$ be a compact, oriented, boundaryless manifold. The global Euler characteristic of the constructible sheaf $\mathcal{F}_A$ equals the intersection product of the characteristic cycle $CC(\mathcal{F}_A)$ with the zero section $T^*_\Omega \Omega$:
\begin{equation}
\chi(\Omega, \mathcal{F}_A) = CC(\mathcal{F}_A) \cdot [T^*_\Omega \Omega] = \sum_{r=0}^n m_r(A) \chi(\Sigma_r(A)),
\end{equation}
where $m_r(A) \in \mathbb{Z}$ are the local Milnor multiplicities of the matrix rank singularity.
\end{theorem}

\begin{proof}
This follows directly from Kashiwara's local index formula for constructible sheaves on smooth manifolds \cite{KashiwaraSchapira1990}, where the characteristic cycle $CC(\mathcal{F}_A) = \sum m_r [\overline{T^*_{\Sigma_r} \Omega}]$ has intersection number with the zero section equal to $\sum m_r \chi(\Sigma_r)$.
\end{proof}

% =========================================================================
% SECTION 4
% =========================================================================
\section{Nonlinear Spectral and Metric Geometry: Weighted \texorpdfstring{$p$}{p}-Laplacians}
\label{sec:nonlinear}

Let $\Omega \subset \mathbb{R}^d$ be a bounded, **connected** domain with Lipschitz boundary $\partial\Omega$, and let $d\mu_A = \Phi(A)(x) \dif x$ with $\Phi(A) \in C^1(\bar{\Omega})$ strictly positive: $\inf_\Omega \Phi(A) \ge c_0 > 0$. For $p \in (1, \infty)$, the weighted $p$-Laplacian operator is defined weakly by
\begin{equation}
\Delta_p^A u \coloneqq \nabla \cdot \left( \Phi(A) \|\nabla u\|^{p-2} \nabla u \right).
\end{equation}

\begin{definition}[\textbf{OBL-007: First $p$-Laplacian Eigenvalue}]
\label{def:p_laplacian}
The first Dirichlet eigenvalue $\lambda_1^{(p)}(A)$ is defined by the Rayleigh--Finsler quotient:
\begin{equation}
\label{eq:rayleigh_p}
\lambda_1^{(p)}(A) \coloneqq \inf_{u \in W_0^{1,p}(\Omega, \mu_A) \setminus \{0\}} \frac{\int_\Omega \|\nabla u\|^p \Phi(A) \dif x}{\int_\Omega |u|^p \Phi(A) \dif x}.
\end{equation}
\end{definition}

\begin{lemma}[\textbf{OBL-007: Positivity and Uniqueness of the Ground State}]
\label{lem:p_laplacian_minimax}
For every $p \in (1, \infty)$, the infimum in \eqref{eq:rayleigh_p} is achieved by a unique positive eigenfunction $u_p \in W_0^{1,p}(\Omega) \cap C^{1,\alpha}(\Omega)$, and $\lambda_1^{(p)}(A) > 0$.
\end{lemma}

\begin{proof}
By Poincaré's inequality on $W_0^{1,p}(\Omega, \mu_A)$, the functional $J(u) = \int_\Omega \|\nabla u\|^p \dif\mu_A$ is coercive and weakly lower semicontinuous on the manifold $S = \{ u \in W_0^{1,p} \mid \int_\Omega |u|^p \dif\mu_A = 1 \}$. Compactness of the Sobolev embedding $W_0^{1,p} \hookrightarrow L^p$ yields an attained minimizer $u_p$. By the strong maximum principle and Harnack inequality for degenerate elliptic equations, connectedness of $\Omega$ ensures that $u_p$ is strictly positive in $\Omega$. Strict convexity of $v \mapsto \|\nabla v\|^p$ ensures uniqueness up to scalar multiplication.
\end{proof}

\begin{theorem}[\textbf{OBL-008: Kawohl--Fridman Cheeger Limit for Functional Realizations}]
\label{thm:cheeger_p_limit}
The normalized first nonlinear eigenvalue family $(\lambda_1^{(p)}(A))^{1/p}$ converges as $p \to \infty$ to the weighted Cheeger isoperimetric constant of the functional realization:
\begin{equation}
\lim_{p \to \infty} \left( \lambda_1^{(p)}(A) \right)^{1/p} = h(A) \coloneqq \inf_{E \subset \Omega} \frac{P(E; \Omega, \Phi(A))}{\mu_A(E)},
\end{equation}
where $P(E; \Omega, \Phi(A)) = \int_{\partial^* E \cap \Omega} \Phi(A) \dif\mathcal{H}^{d-1}$ is the weighted perimeter.
\end{theorem}

\begin{proof}
This asymptotic convergence is established by the foundational theorems of Kawohl and Fridman \cite{KawohlFridman2003} for weighted $p$-Laplace operators, building on the $\infty$-eigenvalue framework of Juutinen, Lindqvist, and Manfredi \cite{JuutinenLindqvistManfredi1999}. 
Specifically, Kawohl and Fridman \cite[Theorem~1.1 and Section~2]{KawohlFridman2003} prove through the viscosity solutions of the Euler--Lagrange equations that for any bounded Lipschitz domain equipped with a strictly positive continuous weight $\Phi(A)$, the variational Rayleigh quotient satisfies $\lim_{p \to \infty} (\lambda_1^{(p)}(A))^{1/p} = \Lambda_\infty(A)$, where $\Lambda_\infty(A)$ is the first eigenvalue of the weighted $\infty$-Laplacian. Furthermore, by linking the weighted coarea formula on level sets of the $\infty$-ground state to the isoperimetric profile, they establish the exact identity $\Lambda_\infty(A) = h(A)$, where $h(A)$ is the weighted Cheeger constant.
\end{proof}

\begin{theorem}[\textbf{OBL-009: Conformal Gromov Hyperbolicity Stability}]
\label{thm:gromov_hyperbolicity}
Let $(\Omega, g_A)$ be a complete, simply connected Riemannian manifold with conformal metric $g_A = \Phi(A)^{2/d} \delta_{ij}$. If the sectional curvature satisfies $\mathrm{Sec}(g_A) \le -\kappa_0 < 0$ uniformly on $\Omega$, then $(\Omega, d_{g_A})$ is complete, geodesic, and Gromov $\delta$-hyperbolic with constant:
\begin{equation}
\delta(A) \le \frac{\ln(1 + \sqrt{2})}{\sqrt{\kappa_0}} < \infty.
\end{equation}
\end{theorem}

\begin{proof}
Since $(\Omega, g_A)$ is a complete, simply connected Riemannian manifold with sectional curvature bounded from above by $-\kappa_0 < 0$, the Cartan--Hadamard theorem ensures that $(\Omega, d_{g_A})$ is a Hadamard manifold, possessing unique minimizing geodesics between any pair of points. By the Rauch comparison theorem and Alexandrov comparison geometry \cite{BridsonHaefliger1999,Gromov1987}, $(\Omega, d_{g_A})$ is a $\mathrm{CAT}(-\kappa_0)$ metric space: every geodesic triangle $\Delta(x, y, z)$ is thinner than a comparison triangle $\bar{\Delta}$ in the model hyperbolic plane $\mathbb{H}^2(-\kappa_0)$. In $\mathbb{H}^2(-\kappa_0)$, the inscribed radius of any ideal triangle is bounded universally by $\kappa_0^{-1/2} \ln(1 + \sqrt{2})$. Consequently, every geodesic triangle in $(\Omega, d_{g_A})$ is $\delta$-slim with $\delta(A) \le \kappa_0^{-1/2} \ln(1 + \sqrt{2})$, proving Gromov hyperbolicity.
\end{proof}

% =========================================================================
% SECTION 5
% =========================================================================
\section{Non-Equilibrium Thermodynamics and Stochastic Dissipation Length}
\label{sec:thermo}

Let $\Phi(A) \in C^2(\Omega, \mathbb{R}_{>0})$ be an unnormalized tensor realization with total mass $Z(A) \coloneqq \int_\Omega \Phi(A)(x) \dif x = \Tr(A) > 0$. Define the potential $V_A(x) \coloneqq -\log \Phi(A)(x)$ and consider the overdamped Langevin diffusion
\begin{equation}
\label{eq:langevin_sde}
\dif X_t = -\nabla V_A(X_t) \dif t + \sqrt{2} \dif W_t.
\end{equation}

\begin{proposition}[\textbf{OBL-010: Well-Posedness and Ergodicity of Langevin Flow}]
\label{prop:langevin_ergodicity}
Assume that the potential $V_A$ is uniformly strongly convex on $\Omega$: $\nabla^2 V_A(x) \ge \kappa I > 0$ for all $x \in \Omega$. Then \eqref{eq:langevin_sde} has a unique stationary probability measure:
\begin{equation}
\dif\mu_A(x) \coloneqq \frac{e^{-V_A(x)}}{Z(A)} \dif x = \frac{\Phi(A)(x)}{\Tr(A)} \dif x,
\end{equation}
and satisfies the Bakry--Émery condition $CD(\kappa, \infty)$, yielding exponential convergence in $\mathcal{W}_2$ with rate $\kappa$.
\end{proposition}

\begin{proof}
The generator $L_A f = \Delta f - \nabla V_A \cdot \nabla f$ is self-adjoint on $L^2(\mu_A)$. Strong convexity $\nabla^2 V_A \ge \kappa I$ implies that the Bakry--Émery curvature-dimension condition $\Gamma_2(f, f) \ge \kappa \Gamma(f, f)$ holds globally, establishing the logarithmic Sobolev inequality with constant $2/\kappa$ via the Bakry--Émery criterion \cite{BakryEmery1985}. By the Otto--Villani theorem \cite{OttoVillani2000}, the logarithmic Sobolev inequality implies the Talagrand transportation inequality $W_2^2(\nu, \mu_A) \le \frac{2}{\kappa} H(\nu | \mu_A)$, which guarantees exponential contraction along the Langevin flow $\mathcal{W}_2(\mu_t, \mu_A) \le e^{-\kappa t} \mathcal{W}_2(\mu_0, \mu_A)$.
\end{proof}

\begin{theorem}[\textbf{OBL-011: Exact Jarzynski Free Energy Identity on Functional Protocols}]
\label{thm:jarzynski_identity}
Let $s \in [0, 1] \mapsto A_s \in \mathcal{S}_+^n$ be a smooth matrix protocol connecting $A_0 = A$ to $A_1 = B$ over duration $\tau$. The stochastic work $W[X_\cdot] = \int_0^\tau \partial_s V_{A_{s(t)}}(X_t) \dot{s}(t) \dif t$ satisfies the exact Jarzynski equality:
\begin{equation}
\mathbb{E}_{\mu_A} \left[ e^{-W} \right] = e^{-\Delta F(A, B)} = \frac{Z(B)}{Z(A)} = \frac{\Tr(B)}{\Tr(A)}.
\end{equation}
\end{theorem}

\begin{proof}
Let $P_F[X_\cdot]$ be the forward path measure of the time-dependent diffusion starting from $\mu_A$, and $P_R[\hat{X}_\cdot]$ the time-reversed path measure starting from $\mu_B$. By Girsanov's theorem and the Crooks fluctuation theorem \cite{Jarzynski1997}, the Radon--Nikodym derivative satisfies $\ln \frac{\dif P_F}{\dif P_R}[X_\cdot] = W[X_\cdot] - \Delta F(A, B)$, where $\Delta F(A, B) = -\ln(Z(B)/Z(A))$. Integrating the identity $P_R[\hat{X}_\cdot] = e^{-(W - \Delta F)} P_F[X_\cdot]$ over the entire path space yields
\begin{equation}
\mathbb{E}_{\mu_A} \left[ e^{-W} \right] = e^{-\Delta F(A, B)} = \frac{Z(B)}{Z(A)} = \frac{\Tr(B)}{\Tr(A)}.
\end{equation}
\end{proof}

\begin{theorem}[\textbf{OBL-012: Finite-Time Thermodynamic Dissipation Bound and Geodesic $\mathcal{W}_2$ Recovery}]
\label{thm:thermo_length_w2}
In the slow-driving limit $\tau \to \infty$, the irreversible dissipated work $\Sigma_{\mathrm{irr}} \coloneqq \langle W \rangle - \Delta F(A, B) \ge 0$ scales inversely with protocol duration $\tau$, and its optimal path defines the thermodynamic length $\mathcal{L}(A, B)$, which satisfies:
\begin{equation}
\lim_{\tau \to \infty} \tau \cdot \inf_\gamma \Sigma_{\mathrm{irr}}[\gamma] = \frac{1}{2} \mathcal{L}^2(A, B) \ge \frac{\kappa}{4} \mathcal{W}_2^2(\mu_A, \mu_B).
\end{equation}
\end{theorem}

\begin{proof}
Expanding the work in powers of $1/\tau$ via linear response theory (Sivak--Crooks \cite{SivakCrooks2012}, Aurell et al. \cite{Aurell2012}), $\langle W \rangle - \Delta F = \frac{1}{\tau} \int_0^1 g_{\mathrm{fric}}(A_s)(\dot{A}_s, \dot{A}_s) \dif s + \mathcal{O}(\tau^{-2})$, where $g_{\mathrm{fric}}$ is the friction metric computed from the time-integrated auto-correlation of conjugate forces along the Langevin diffusion. By the Poincaré inequality with spectral gap $\kappa$ of the generator $L_A$ (guaranteed by $\nabla^2 V_A \ge \kappa I$), the friction tensor dominates the Otto Wasserstein metric: $g_{\mathrm{fric}} \ge \frac{\kappa}{2} g_{\mathrm{Otto}}$. Minimizing over paths gives the geodesic inequality $\mathcal{L}^2(A, B) \ge \frac{\kappa}{2} \mathcal{W}_2^2(\mu_A, \mu_B)$, yielding the bound $\frac{1}{2}\mathcal{L}^2 \ge \frac{\kappa}{4}\mathcal{W}_2^2$.
\end{proof}

% =========================================================================
% SECTION 6
% =========================================================================
\section{Bipartite Kernel Operators and Holographic Entanglement Duals}
\label{sec:bipartite}

We now perform the decisive codomain leap of Volume~III: the realization of a positive semi-definite matrix $A \in \mathcal{S}_+^n$ on a composite bipartite system whose base geometry is the Cartesian product $\Omega \coloneqq \Omega_1 \times \Omega_2$ of two Riemannian manifolds $\Omega_1$ and $\Omega_2$. The continuous state space is the tensor product Hilbert space:
\begin{equation}
\mathcal{H} \coloneqq L^2(\Omega_1 \times \Omega_2) \cong L^2(\Omega_1) \otimes L^2(\Omega_2).
\end{equation}
Let $\{\psi_k\}_{k=1}^n$ be a fixed, canonical orthonormal spatial frame in $L^2(\Omega)$ (e.g. Laplacian eigenfunctions or localized wavelet packets). For any matrix $A = (A_{jk}) \in \mathcal{S}_+^n$, its continuous realization is defined algebraically by:
\begin{equation}
\label{eq:kernel_alg_def}
K_A(x, y) \coloneqq \sum_{j, k=1}^n A_{jk} \psi_j(x) \overline{\psi_k(y)}, \qquad x, y \in \Omega_1 \times \Omega_2.
\end{equation}

\begin{proposition}[\textbf{OBL-013: Positive Kernel Codomain Realization and Constructive Diagonal Recovery}]
\label{prop:kernel_codomain_psd}
Let $A \in \mathcal{S}_+^n$ have unitary spectral decomposition $A = U \operatorname{diag}(\lambda_1, \dots, \lambda_n) U^\dagger$. Define the rotated functional basis:
\begin{equation}
\phi_j(x) \coloneqq \sum_{k=1}^n U_{kj} \psi_k(x), \qquad j = 1, \dots, n.
\end{equation}
Then $\{\phi_j\}_{j=1}^n$ is an orthonormal family in $L^2(\Omega_1 \times \Omega_2)$, the kernel $K_A$ in \eqref{eq:kernel_alg_def} is identically diagonalized:
\begin{equation}
K_A(x, y) = \sum_{j=1}^n \lambda_j \phi_j(x) \overline{\phi_j(y)},
\end{equation}
and the induced integral operator $T_A$ is self-adjoint, trace-class, positive semi-definite with $\Tr(T_A) = \Tr(A)$, whose diagonal restriction constructively coincides with the scalar realization $\Phi_{\mathrm{scalar}}(A)(x) \coloneqq \sum_{j,k} A_{jk} \psi_j(x) \overline{\psi_k(x)}$.
\end{proposition}

\begin{proof}
Orthonormality of $\{\phi_j\}$ is guaranteed by the unitarity of $U$:
\begin{equation}
\langle \phi_i, \phi_j \rangle_{L^2} = \sum_{k, l=1}^n U_{ki}^* U_{lj} \langle \psi_k, \psi_l \rangle_{L^2} = \sum_{k=1}^n U_{ki}^* U_{kj} = (U^\dagger U)_{ij} = \delta_{ij}.
\end{equation}
Expanding $K_A(x, y)$ in terms of $\lambda_j$ and $\phi_j$, we compute $\langle f, T_A f \rangle = \sum \lambda_j |\langle f, \phi_j \rangle|^2 \ge 0$ since $\lambda_j \ge 0$. The trace is $\int_\Omega K_A(x, x) \dif x = \sum \lambda_j \|\phi_j\|^2 = \sum \lambda_j = \Tr(A)$. Finally, $K_A(x, x) = \sum \lambda_j |\phi_j(x)|^2 = \sum_{j,k} A_{jk} \psi_j(x) \overline{\psi_k(x)} = \Phi_{\mathrm{scalar}}(A)(x)$ holds identically by matrix multiplication.
\end{proof}

\begin{theorem}[\textbf{OBL-014: Continuous Partial Trace and Support-Restricted Modular Hamiltonian}]
\label{thm:modular_hamiltonian_spec}
On the bipartite product space $\Omega = \Omega_1 \times \Omega_2$, the continuous partial trace of $K_A$ over subsystem $\Omega_2$ is defined by:
\begin{equation}
\rho_{\Omega_1}(x_1, y_1) \coloneqq \int_{\Omega_2} K_A(x_1, z_2; y_1, z_2) \dif z_2.
\end{equation}
The partial trace $\rho_{\Omega_1}$ is a self-adjoint, positive semi-definite trace-class operator on $L^2(\Omega_1)$ of finite rank $r \le n$. Restricted to its active support subspace $\mathcal{H}_{\mathrm{supp}}^{(1)} \coloneqq (\ker \rho_{\Omega_1})^\perp \subset L^2(\Omega_1)$ of dimension $r$, $\rho_{\Omega_1}$ is strictly positive, and the support-restricted modular Hamiltonian
\begin{equation}
K_{\Omega_1} \coloneqq -\log\left(\left.\rho_{\Omega_1}\right|_{\mathcal{H}_{\mathrm{supp}}^{(1)}}\right)
\end{equation}
is a strictly positive, self-adjoint operator with discrete finite spectrum $\{E_k = -\log p_k\}_{k=1}^r$. Under continuous Mercer regularization $K_A^\epsilon = K_A + \epsilon (k_{\Omega_1} \otimes k_{\Omega_2})$, where $k_{\Omega_i}$ are strictly positive definite reproducing kernels, $\ker(\rho_{\Omega_1}^\epsilon) = \{0\}$ and $K_{\Omega_1}^\epsilon$ has an infinite discrete spectrum accumulating at $+\infty$.
\end{theorem}

\begin{proof}
Writing $K_A = \sum_{j=1}^n \lambda_j |v_j\rangle\langle v_j|$ with $v_j \in L^2(\Omega_1 \times \Omega_2)$, the partial trace is $\rho_{\Omega_1} = \sum_{j=1}^n \lambda_j \Tr_2(|v_j\rangle\langle v_j|)$, which is a positive semi-definite operator of rank $r \le n$ on $L^2(\Omega_1)$. On the finite-dimensional invariant subspace $\mathcal{H}_{\mathrm{supp}}^{(1)} = \im(\rho_{\Omega_1})$, $\rho_{\Omega_1}$ has non-zero eigenvalues $p_1 \ge \dots \ge p_r > 0$. Thus $K_{\Omega_1} = \sum_{k=1}^r (-\log p_k) |e_k\rangle\langle e_k|$ is self-adjoint and positive on $\mathcal{H}_{\mathrm{supp}}^{(1)}$. For the regularized kernel $K_A^\epsilon$, strict positive definiteness of $k_{\Omega_1}$ ensures trivial kernel on $L^2(\Omega_1)$, with compact resolvent yielding discrete spectrum $E_k^\epsilon \to +\infty$.
\end{proof}

\begin{theorem}[\textbf{OBL-015: Canonical Purification and Reflected Entropy Monotonicity}]
\label{thm:reflected_entropy_prop}
Let $\rho_{12}$ be the normalized density operator of $K_A$ on $\mathcal{H}_{\mathrm{supp}}^{(1)} \otimes \mathcal{H}_{\mathrm{supp}}^{(2)}$, and let $|\sqrt{\rho_{12}}\rangle \in (\mathcal{H}_{\mathrm{supp}}^{(1)} \otimes \mathcal{H}_{\mathrm{supp}}^{(2)}) \otimes (\mathcal{H}_{\mathrm{supp}}^{(1^*)} \otimes \mathcal{H}_{\mathrm{supp}}^{(2^*)})$ be its canonical GNS purification. The reflected entropy
\begin{equation}
S_R(\Omega_1 : \Omega_2) \coloneqq -\Tr\left( \sigma_{1 1^*} \log \sigma_{1 1^*} \right), \qquad \sigma_{1 1^*} \coloneqq \Tr_{2 2^*}\left( |\sqrt{\rho_{12}}\rangle \langle\sqrt{\rho_{12}}| \right)
\end{equation}
is non-negative, symmetric $S_R(\Omega_1 : \Omega_2) = S_R(\Omega_2 : \Omega_1)$, and satisfies the universal lower bound:
\begin{equation}
S_R(\Omega_1 : \Omega_2) \ge I(\Omega_1 : \Omega_2) \coloneqq S(\Omega_1) + S(\Omega_2) - S(\Omega_1 \Omega_2).
\end{equation}
\end{theorem}

\begin{proof}
Positivity and symmetry follow from properties of canonical purifications in Tomita--Takesaki theory. The lower bound by mutual information is established via the monotonicity of relative entropy under partial tracing of the purifying auxiliary systems $1^*, 2^*$, as proven by Dutta and Faulkner \cite{DuttaFaulkner2019}.
\end{proof}

\begin{theorem}[\textbf{OBL-016: Holographic Entanglement Wedge Cross-Section Inequality}]
\label{thm:holographic_ew_inequality}
In any holographic geometry where the functional bulk geometry $(M, g_A)$ satisfies the Einstein equations with null energy condition and boundary dual density operator $\rho_A$, the reflected entropy satisfies:
\begin{equation}
S_R(\Omega_1 : \Omega_2) \ge 2 E_W(A; \Omega_1, \Omega_2) = \frac{1}{2 G_N} \min_{\Sigma_{12}} \mathrm{Area}_{g_A}(\Sigma_{12}),
\end{equation}
where $\Sigma_{12}$ bisects the entanglement wedge of $\Omega_1 \cup \Omega_2$.
\end{theorem}

\begin{proof}
This follows from the gravitational replica trick for reflected entropy (Dutta--Faulkner \cite{DuttaFaulkner2019}), where the replicated boundary manifold contracts in the bulk to a conical cosmic brane defect whose minimal bounding area precisely equals twice the entanglement wedge cross section $E_W$.
\end{proof}

% =========================================================================
% SECTION 7
% =========================================================================
\section{Simplicial Pascal Simplex and Barnes--Kigami Residue Spectrum}
\label{sec:simplicial}

Let $\Delta_m \coloneqq \{ u \in \mathbb{R}_+^m \mid \sum_{i=1}^m u_i = 1 \}$ be the standard continuous Pascal simplex, equipped with the Dirichlet probability measure $\dif\mu_\alpha(u) = \frac{\Gamma(\sum \alpha_i)}{\prod \Gamma(\alpha_i)} \prod u_i^{\alpha_i - 1} \dif u$.

\begin{proposition}[\textbf{OBL-017: Simplicial Mellin Transform Holomorphy}]
\label{prop:simplicial_mellin_conv}
For any strictly positive functional realization $\Phi(A) \in C^1(\Delta_m)$, the simplicial Mellin transform
\begin{equation}
\mathcal{Z}_A(s) \coloneqq \int_{\Delta_m} \left( \Phi(A)(u) \right)^s \dif\mu_\alpha(u)
\end{equation}
is a holomorphic function of $s$ in the right half-plane $\operatorname{Re}(s) > -\min_i \alpha_i$.
\end{proposition}

\begin{proof}
Because $\Phi(A)$ is bounded and strictly positive on the compact simplex $\Delta_m$, $0 < c_1 \le \Phi(A)(u) \le c_2 < \infty$. Thus $|(\Phi(A)(u))^s| \le c_2^{\operatorname{Re}(s)}$ for $\operatorname{Re}(s) \ge 0$, and the Dirichlet density is integrable for $\operatorname{Re}(s) > -\min \alpha_i$. Differentiating under the integral sign proves complex analyticity in this half-plane.
\end{proof}

\begin{theorem}[\textbf{OBL-018: Meromorphic Continuation and Barnes Residue Spectrum}]
\label{thm:barnes_meromorphic_continuation}
The partition function $\mathcal{Z}_A(s)$ extends meromorphically to the entire complex plane $\mathbb{C}$, with isolated simple poles located at $s_k = -\alpha_j - k$ for $k \in \mathbb{N}_0$, whose residues define the discrete Barnes spectrum $\{\operatorname{Res}_{s = s_k} \mathcal{Z}_A(s)\}_{k \in \mathbb{N}}$.
\end{theorem}

\begin{proof}
Expanding $\Phi(A)(u)$ into its Taylor series around the vertices $e_j \in \Delta_m$, the singular contributions near vertex $e_j$ reduce to multivariable Beta integrals whose numerator contains Euler gamma functions $\Gamma(s + \alpha_j + k)$. These gamma functions possess simple poles at $s + \alpha_j + k = 0, -1, -2, \dots$ with non-vanishing residues, yielding the discrete Barnes residue spectrum.
\end{proof}

\begin{corollary}[\textbf{OBL-019: Kigami Fractal Spectral Dimension Identity}]
\label{cor:kigami_spectral_dim}
For a self-similar fractal realization generated by iterated simplex decimations with $N$ cells and harmonic resistance scaling factor $\rho \in (0, 1)$, the dominant pole $s_0$ of the renormalized Mellin transform gives the exact Kigami spectral dimension:
\begin{equation}
s_0 = -\frac{d_s}{2} = -\frac{\log N}{\log(N / \rho)} \implies d_s = \frac{2 \log N}{\log(N / \rho)}.
\end{equation}
\end{corollary}

\begin{proof}
Under Kigami's decimation of the resistance form on the fractal simplex \cite{Kigami2001}, the heat kernel trace scales asymptotically as $t^{-d_s/2}$ where $d_s = \frac{2\log N}{\log(N/\rho)}$. Taking the Mellin transform links the heat trace directly to the partition function poles, placing the dominant pole at $s_0 = -d_s/2$. For the standard benchmark of the Sierpinski gasket ($N=3, \rho = 3/5$), this yields $N/\rho = 3/(3/5) = 5$, hence $d_s = \frac{2\log 3}{\log 5} \approx 1.3652$, in exact agreement with Kigami's theorem.
\end{proof}

% =========================================================================
% SECTION 8
% =========================================================================
\section{Federer Reach and Medial Axis Geometry of Level Surfaces}
\label{sec:federer}

\begin{theorem}[\textbf{OBL-020: Level Set Federer Reach Lower Bound}]
\label{thm:federer_reach_bound}
Let $\Phi(A) \in C^2(\Omega, \mathbb{R})$. For any regular value $t$ with $\epsilon_0 \coloneqq \inf_{\Sigma_t} \|\nabla \Phi(A)\| > 0$ and $M \coloneqq \sup_{\Sigma_t} \|\nabla^2 \Phi(A)\|_{\mathrm{op}} < \infty$, let $d_{\mathrm{sep}}(\Sigma_t) \coloneqq \inf \{ \|x - y\| \mid x, y \in \Sigma_t, \nu(x) = -\nu(y) \} > 0$ be the non-local bottleneck separation distance. Then the Federer reach of $\Sigma_t = \Phi(A)^{-1}(t)$ satisfies:
\begin{equation}
\mathrm{reach}(\Sigma_t) \ge \min\left( \frac{\epsilon_0}{M}, \frac{1}{2} d_{\mathrm{sep}}(\Sigma_t) \right) > 0.
\end{equation}
When the sublevel set $K_t(A)$ is strictly convex, $d_{\mathrm{sep}}(\Sigma_t) = \infty$, yielding the pure curvature bound $\mathrm{reach}(\Sigma_t) \ge \frac{\epsilon_0}{M}$.
\end{theorem}

\begin{proof}
The second fundamental form $\mathrm{II}$ of the regular level set $\Sigma_t$ satisfies:
\begin{equation}
\|\mathrm{II}(v, v)\| \le \frac{\|\nabla^2 \Phi(A)\|_{\mathrm{op}}}{\|\nabla \Phi(A)\|} \le \frac{M}{\epsilon_0}
\end{equation}
for any unit tangent vector $v \in T_x \Sigma_t$. By Federer's theorem \cite{Federer1959}, the reach of a compact $C^2$ submanifold is determined by the minimum between the local radius of curvature $\kappa_{\max}^{-1} \ge \epsilon_0 / M$ and half the non-local bottleneck distance $\frac{1}{2} d_{\mathrm{sep}}(\Sigma_t)$. For strictly convex level sets, no two distinct points share anti-parallel normals with connecting segment, so $d_{\mathrm{sep}} = \infty$, reducing the bound to $\epsilon_0 / M$.
\end{proof}

\begin{theorem}[\textbf{OBL-021: Medial Axis Avoidance and Weyl--Federer Tube Volume Expansion}]
\label{thm:tube_volume_reach}
Let $\Sigma_t$ have reach $R_t > 0$ established by Theorem~\ref{thm:federer_reach_bound}. For any $r < R_t$, the tubular neighborhood $U_r(\Sigma_t)$ does not intersect the medial axis $\mathcal{M}(\Phi(A))$, the metric projection $\Proj_{\Sigma_t} \colon U_r \to \Sigma_t$ is single-valued and Lipschitz, and the tube volume satisfies the exact polynomial expansion:
\begin{equation}
\mathrm{Vol}_d(U_r(\Sigma_t)) = \sum_{k=0}^{\lfloor \frac{d-1}{2} \rfloor} \frac{2 r^{2k+1}}{2k+1} \int_{\Sigma_t} H_{2k}(x) \dif\mathcal{H}^{d-1}(x),
\end{equation}
where $H_{2k}$ are the even elementary symmetric functions of principal curvatures.
\end{theorem}

\begin{proof}
By definition of reach, the set of points with non-unique projection $\Unp(\Sigma_t)^c$ has distance at least $R_t$ from $\Sigma_t$. Thus for $r < R_t$, $U_r \cap \mathcal{M} = \emptyset$. The volume of the tubular neighborhood is obtained by integrating the Jacobian of the normal exponential map, which by Weyl's tube formula expands into integrated curvature forms $H_{2k}$ \cite{Federer1959}.
\end{proof}

% =========================================================================
% SECTION 9
% =========================================================================
\section{Grand Invariant Taxonomy and Synthesis}
\label{sec:taxonomy}

Table~\ref{tab:taxonomy} summarizes the evolutionary taxonomy of geometric and topological measures on functional realizations across the entire trilogy.

\begin{table}[htbp]
\centering
\small
\caption{Grand Taxonomy of Invariants across Volumes I, II, and III.}
\label{tab:taxonomy}
\begin{tabular}{lllll}
\toprule
\textbf{Invariant} & \textbf{Codomain} & \textbf{Geometric Structure} & \textbf{Symmetry Group} & \textbf{Volume} \\
\midrule
$\lambda_j(A), \Tr(A)$ & $\mathbb{R}^n$ & Euclidean / Linear & $O(n) / U(n)$ & Classical \\
$\|\Phi(A)\|_{W^{1,p}}$ & $\mathbb{R}$ & Sobolev / BV & Diffeomorphisms & Vol. I \\
$\mathcal{W}_2(\mu_A, \mu_B)$ & $\mathbb{R}_+$ & Optimal Transport & Diff-preserving & Vol. II \\
$\kappa_{\mathrm{BE}}(A)$ & $\mathbb{R}$ & Bakry--Émery Curvature & Conformal & Vol. II \\
$\mathrm{Dgm}_k(A)$ & $\mathbb{R}$ & Persistent Homology & Homeomorphisms & Vol. II \\
$\mathcal{W}(\Phi(A))$ & $\mathbb{R}$ & Willmore Bending & Conformal & Vol. II \\
\midrule
$c_k^{\mathrm{EH}}(A, t), c_{\mathrm{HZ}}$ & $\mathbb{R}$ & Symplectic Capacity & $\mathrm{Symp}(\mathbb{R}^{2n}, \omega_0)$ & \textbf{Vol. III} \\
$CC(\mathcal{F}_A)$ & $\mathrm{Sym}_n(\mathbb{R})$ & Microlocal Sheaves & Stratified Diffeo & \textbf{Vol. III} \\
$\lambda_1^{(p)}(A)$ & $\mathbb{R}$ & Nonlinear Finsler & Monotone & \textbf{Vol. III} \\
$\delta(A)$ & $\mathbb{R}$ & Gromov Hyperbolic & Coarse Quasi-isom & \textbf{Vol. III} \\
$\mathcal{L}(A, B)$ & $\mathbb{R}_+$ & Non-Equilibrium Thermo & Gauge Flow & \textbf{Vol. III} \\
$K_{\Omega_1}, S_R, E_W$ & $L^2(\Omega \times \Omega)$ & Holographic Entanglement & Unitary Bipartite & \textbf{Vol. III} \\
$\mathcal{Z}_A(s), s_0$ & $\Delta_m$ & Barnes / Kigami Fractal & Simplicial Scale & \textbf{Vol. III} \\
$\mathrm{reach}(\Sigma_t)$ & $\mathbb{R}$ & Federer Medial Axis & Isometric Tube & \textbf{Vol. III} \\
\bottomrule
\end{tabular}
\end{table}

% =========================================================================
% REFERENCES
% =========================================================================
\begin{thebibliography}{99}

\bibitem{Aurell2012}
E.~Aurell, C.~Mej\'ia-Monasterio, and P.~Muratore-Ginanneschi,
\emph{Optimal protocols and optimal transport in stochastic thermodynamics},
Phys. Rev. Lett. \textbf{106} (2011), 250601.

\bibitem{BakryEmery1985}
D.~Bakry and M.~\'Emery,
\emph{Diffusions hypercontractives},
S\'eminaire de probabilit\'es de Strasbourg \textbf{19} (1985), 177--206.

\bibitem{DuttaFaulkner2019}
K.~Dutta and T.~Faulkner,
\emph{A canonical purification for the entanglement wedge cross-section},
J. High Energ. Phys. \textbf{2021} (2021), 28.

\bibitem{EkelandHofer1989}
I.~Ekeland and H.~Hofer,
\emph{Symplectic topology and Hamiltonian dynamics I},
Math. Z. \textbf{200} (1989), 355--378.

\bibitem{Federer1959}
H.~Federer,
\emph{Curvature measures},
Trans. Amer. Math. Soc. \textbf{93} (1959), 418--491.

\bibitem{Gromov1987}
M.~Gromov,
\emph{Hyperbolic groups},
Essays in group theory, Springer, New York, 1987, pp.~75--263.

\bibitem{HoferZehnder1994}
H.~Hofer and E.~Zehnder,
\emph{Symplectic Invariants and Hamiltonian Dynamics},
Birkh\"auser Basel, 1994.

\bibitem{Jarzynski1997}
C.~Jarzynski,
\emph{Nonequilibrium equality for free energy differences},
Phys. Rev. Lett. \textbf{78} (1997), 2690.

\bibitem{KashiwaraSchapira1990}
M.~Kashiwara and P.~Schapira,
\emph{Sheaves on Manifolds},
Grundlehren der mathematischen Wissenschaften, vol.~292, Springer-Verlag, Berlin, 1990.

\bibitem{BridsonHaefliger1999}
M.~R. Bridson and A.~Haefliger,
\emph{Metric Spaces of Non-Positive Curvature},
Grundlehren der mathematischen Wissenschaften, vol.~319, Springer, Berlin, 1999.

\bibitem{JuutinenLindqvistManfredi1999}
P.~Juutinen, P.~Lindqvist, and J.~J. Manfredi,
\emph{The $\infty$-eigenvalue problem},
Arch. Ration. Mech. Anal. \textbf{148} (1999), 89--105.

\bibitem{OttoVillani2000}
F.~Otto and C.~Villani,
\emph{Generalization of an inequality by Talagrand and links with the logarithmic Sobolev inequality},
J. Funct. Anal. \textbf{173} (2000), 361--400.

\bibitem{KawohlFridman2003}
B.~Kawohl and V.~Fridman,
\emph{Isoperimetric estimates for the first eigenvalue of the $p$-Laplace operator and the Cheeger constant},
Nonlinear Anal. \textbf{55} (2003), 107--118.

\bibitem{Kigami2001}
J.~Kigami,
\emph{Analysis on Fractals},
Cambridge Tracts in Mathematics, Cambridge University Press, 2001.

\bibitem{SilvaFilho2026a}
R.~M. Silva-Filho,
\emph{A Functorial Bridge from Continuous Tensor Manifolds to 4-Dimensional Spacetime Cobordisms},
Zenodo (2026), DOI: \href{https://doi.org/10.5281/zenodo.22441676}{10.5281/zenodo.22441676}.

\bibitem{SilvaFilho2026b}
R.~M. Silva-Filho,
\emph{Unified Quantum Gravity and Continuous Tensor Field Theory},
Zenodo (2026), DOI: \href{https://doi.org/10.5281/zenodo.22290043}{10.5281/zenodo.22290043}.

\bibitem{SilvaFilho2026c}
R.~M. Silva-Filho,
\emph{Advanced Geometric Measures and Topological Invariants on Functional Realizations of Matrices and Higher-Order Tensors (Beyond the Spectrum II)},
Preprint (2026).

\bibitem{SivakCrooks2012}
D.~A. Sivak and G.~E. Crooks,
\emph{Thermodynamic metrics and optimal paths},
Phys. Rev. Lett. \textbf{108} (2012), 190602.

\bibitem{Viterbo1992}
C.~Viterbo,
\emph{Symplectic topology as the geometry of generating functions},
Math. Ann. \textbf{292} (1992), 685--710.

\end{thebibliography}

\end{document}
